% =====================================================================
%  demo.tex --- the Lucid beamer theme, explained by example
%
%  This deck IS the manual. Each slide shows a component and, under it,
%  the exact code that produced it. Copy a block off the slide, or open
%  the matching file in snippets/ and copy it from there.
%
%  Build:  pdflatex demo.tex     (run it twice -- the outline slides
%                                 need a second pass to fill in)
%
%  Setup for Overleaf and for VS Code + LaTeX Workshop is in README.md.
%
%  Why snippets/ exists: a beamer frame marked [fragile] ends at the
%  first \end{frame} on a line, including one that is only sample text
%  inside a listing. Reading the samples from files sidesteps that
%  entirely, and gives you files you can copy straight out.
% =====================================================================

\documentclass[11pt,aspectratio=43]{beamer}
\usetheme{lucid}

\usepackage{tikz}
\usepackage{listings}

% Code blocks. Needed only by THIS deck, because this deck prints its
% own source. A normal deck does not need listings.
\lstset{
  language=[LaTeX]TeX,
  basicstyle=\ttfamily\fontsize{6.8}{8.2}\selectfont\color{lucidInk},
  keywordstyle=\color{lucidAccent},
  commentstyle=\color{lucidMuted}\itshape,
  backgroundcolor=\color{lucidBand!40},
  breaklines=true,
  breakindent=8pt,
  columns=fullflexible,
  keepspaces=true,
  frame=none,
  xleftmargin=5pt,
  xrightmargin=3pt,
  aboveskip=4pt,
  belowskip=2pt,
  showstringspaces=false,
}
\newcommand{\howto}[1]{\lstinputlisting{snippets/#1.tex}}

\title{The Lucid Beamer Theme}
\subtitle{Every slide shows the code that made it}
\lucidevent{A copy-and-paste manual}
\date{}

% Identity is off by default. Uncomment to switch it on.
% \author{Your Name}
% \institute{Your Institution}
% \lucidlogo{logo.pdf}

\begin{document}

% --- Title slide -----------------------------------------------------
\begin{frame}[plain]
  \titlepage
\end{frame}

\section{Getting started}

\begin{frame}{Start a deck: two lines}
  \begin{itemize}
    \item Put \texttt{beamerthemelucid.sty} next to your \texttt{.tex} file
  \end{itemize}
  \howto{start}
  \begin{itemize}
    \item Use \texttt{aspectratio=169} for widescreen
  \end{itemize}
  \citeline{Compile twice. The outline slides need a second pass. \quad snippets/start.tex}
\end{frame}

\begin{frame}{The title slide}
  \begin{itemize}
    \item \term{No name, institution or logo appears unless you set one}
    \item \texttt{[plain]} drops the footer, which a title slide should not have
  \end{itemize}
  \howto{title-slide}
  \citeline{snippets/title-slide.tex}
\end{frame}

\section{Text}

\begin{frame}{Bullets: two levels, and a third to avoid}
  \begin{itemize}
    \item Level one
      \begin{itemize}
        \item Level two
          \begin{itemize}
            \item Level three is muted on purpose
          \end{itemize}
      \end{itemize}
  \end{itemize}
  \howto{bullets}
  \citeline{Aim for two top-level bullets and 25--30 words a slide. \quad snippets/bullets.tex}
\end{frame}

\begin{frame}{Emphasis: three devices, and none of them is bold}
  \begin{itemize}
    \item A \term{primary term}, a \altterm{secondary term}, something \hl{highlighted}
    \item Each one survives greyscale printing on its own
  \end{itemize}
  \howto{emphasis}
  \citeline{snippets/emphasis.tex}
\end{frame}

\section{Figures}

\begin{frame}{A figure, uncaptioned}
  \begin{center}
    \begin{tikzpicture}[scale=0.38]
      \fill[lucidBand] (0,0) rectangle (9,4);
      \foreach \x/\h in {0.7/3.1, 2.2/2.6, 3.7/3.3, 5.2/2.0, 6.7/2.8}
        {\fill[lucidAccent] (\x,0.4) rectangle (\x+0.9,\h);}
      \draw[lucidGold,line width=1.2pt] (0.3,2.2) -- (8.7,2.2);
    \end{tikzpicture}
  \end{center}
  \howto{figure}
  \citeline{snippets/figure.tex}
\end{frame}

\begin{frame}{A figure beside text}
  \begin{figuretext}[0.55]
    \begin{itemize}
      \item The left pane holds your points
      \item The argument sets its share
    \end{itemize}
    \nextpane
    \begin{tikzpicture}[scale=0.30]
      \draw[lucidMuted,line width=0.4pt] (0,0) rectangle (10,7.5);
      \fill[lucidBand] (0,6.3) rectangle (10,7.5);
      \fill[lucidGold] (0,6.18) rectangle (10,6.32);
      \draw[lucidAccent,dashed,line width=0.3pt] (0.6,0) -- (0.6,7.5);
    \end{tikzpicture}
  \end{figuretext}
  \howto{figure-text}
  \citeline{snippets/figure-text.tex}
\end{frame}

\section{Tables, maths, closings}

\begin{frame}{A table}
  \begin{lucidtable}
    \begin{tabular}{lr}
      \toprule
      \thead{Method & Score} \\
      \midrule
      Baseline & 62 \\
      Ours     & \term{78} \\
      \bottomrule
    \end{tabular}
  \end{lucidtable}
  \howto{table}
  \citeline{Horizontal rules only. \quad snippets/table.tex}
\end{frame}

\begin{frame}{An equation}
  \begin{itemize}
    \item Equations are unnumbered, and sit inside the bullet that introduces them
      \[ q_j \;=\; \frac{e^{\lambda\,\mathrm{EU}_j(x)}}{\sum_i e^{\lambda\,\mathrm{EU}_i(x)}} \]
  \end{itemize}
  \howto{equation}
  \citeline{snippets/equation.tex}
\end{frame}

\begin{frame}{A source line, and a closing statement}
  \takeaway{The one sentence you want them to leave with.}
  \howto{closing}
  \citeline{A source line sits just above the footer rule, like this one.}
\end{frame}

\begin{frame}{Sections make their own divider}
  \begin{itemize}
    \item Every \texttt{section} inserts an outline slide with the current section marked
    \item You have passed four of them already
  \end{itemize}
  \howto{sections}
  \citeline{snippets/sections.tex}
\end{frame}

\section{Options and backup}

\begin{frame}{Theme options}
  \begin{lucidtable}
    \begin{tabular}{@{}ll@{}}
      \toprule
      \thead{Option & Values, default first} \\
      \midrule
      \texttt{palette}    & \texttt{green}, \texttt{maroon}, \texttt{slate} \\
      \texttt{font}       & \texttt{cabin}, \texttt{helvet}, \texttt{none} \\
      \texttt{footer}     & \texttt{full}, \texttt{page}, \texttt{none} \\
      \texttt{pagenumber} & \texttt{plain}, \texttt{total} \\
      \texttt{outline}    & \texttt{auto}, \texttt{manual} \\
      \texttt{align}      & \texttt{top}, \texttt{center} \\
      \texttt{blind}      & \texttt{false}, \texttt{true} \\
      \bottomrule
    \end{tabular}
  \end{lucidtable}
  \howto{options}
\end{frame}

\begin{frame}{Anonymous submission}
  \begin{itemize}
    \item \term{Nothing identifying is emitted unless you set it}
    \item \texttt{blind} suppresses author, institution and logo even when set, so a
      deck flips to anonymous without editing its front matter
  \end{itemize}
  \howto{blind}
  \citeline{snippets/blind.tex}
\end{frame}

\begin{frame}{Backup slides}
  \begin{itemize}
    \item \texttt{backupframes} stops later slides inflating the page count
  \end{itemize}
  \howto{backup}
  \citeline{The next slide is a backup slide. Watch the page number restart.}
\end{frame}

\backupframes

\begin{frame}{Backup: the numbering restarted}
  \begin{itemize}
    \item This is a backup slide
    \item The footer count started again at 1
  \end{itemize}
\end{frame}

\end{document}
